\chapter{System Context}
\ChapterLabel{SystemContext}
\index{system boundary}
\index{analog input}
\index{digital readings}

\section{Chapter Responsibility}
\label{Section:SystemContextResponsibility}
Define the ADC system boundary, external participants, exchanged signals, and
unresolved operating assumptions.

\noindent\textbf{Expected result:} A shared boundary from which measurable
requirements and interface decisions can be developed.

\section{Problem and Intended Outcome}
The project owner needs an embedded system that accepts an analog input and
converts it to sampled numeric ADC readings. The intended demonstration makes
those readings observable for comparison with a known input. The demonstration
method and pass criteria are proposals until the course reviewer confirms the
decisions in \hyperref[Table:WeekOneTBDs]{the Week 1 decision register}.

\section{System Boundary}
The proposed system comprises the selected embedded target, ADC capability,
firmware needed to initiate acquisition and obtain readings, and an output
mechanism for observing those readings. Whether the ADC is integrated into
the processor or connected externally is TBD. Any necessary input protection,
conditioning, reference, or filtering must be allocated during hardware review;
none is assumed to be present on an unselected board.

The analog source and its measured physical quantity are outside the boundary.
The external observer or receiving tool is outside; the target-side mechanism
that exposes readings is inside. Power and programming equipment are external
dependencies. Figure~\ref{Figure:ADCContext} presents this proposed boundary
without selecting a board, communication protocol, or signal range.

\begin{figure}[htbp]
\centering
\begin{tikzpicture}[
 box/.style={draw=dsnBlue,fill=dsnLightBlue,rounded corners,align=center,
 text width=4cm,minimum height=1.1cm,font=\small},
 ext/.style={draw=dsnGray,align=center,text width=2.8cm,
 minimum height=1cm,font=\small},
 >=Stealth]
\node[box,minimum height=2.6cm] (sys) at (0,0)
 {Embedded ADC system\\Target, ADC, firmware\\and output mechanism\\Hardware TBD};
\node[ext] (src) at (-5.5,0) {Analog source};
\node[ext] (obs) at (5.5,0) {Observer /\\receiving tool};
\node[ext] (power) at (0,3.2) {External power};
\node[ext] (tools) at (0,-3.2) {Programming /\\debug tools};
\draw[->] (src.east) -- node[above,font=\scriptsize]{Analog input} (sys.west);
\draw[->] (sys.east) -- node[above,font=\scriptsize]{ADC readings} (obs.west);
\draw[->] (power.south) -- node[right,font=\scriptsize]{Power} (sys.north);
\draw[<->] (tools.north) -- node[right,font=\scriptsize]{Image / diagnostics} (sys.south);
\end{tikzpicture}

\caption{Proposed ADC system context; electrical and digital interface details are TBD.}
\label{Figure:ADCContext}
\end{figure}

\section{External Participants and Interfaces}
\small
\begin{tabularx}{\linewidth}{L{2.9cm}Y Y}
\caption{External Interactions\label{Table:ADCExternalInteractions}}\\
\toprule
Participant & Exchange at the boundary & Ownership and unresolved contract\\
\midrule
\endfirsthead
\toprule
Participant & Exchange at the boundary & Ownership and unresolved contract\\
\midrule
\endhead
\bottomrule
\endlastfoot
Analog source & Analog input into the system; electrical reference or return
as required by the chosen interface. &
Jack Marino identifies the source, range, impedance, bandwidth, reference,
protection, and calibration basis under TBD-02.\\
Observer or receiving tool & Sampled numeric readings from the system;
configuration or acquisition commands only if later selected. &
Jack Marino selects the output interface, representation, scaling, framing,
and observation tool under TBD-04.\\
Power source & Electrical power into the system. &
Jack Marino establishes supply limits, grounding, startup behavior, and
available current under TBD-01 and TBD-05.\\
Development and programming tools & Firmware image and programming/debug
operations to the target; diagnostic information back to the tools. &
Jack Marino selects the toolchain, host environment, programming connection,
and compatible versions under TBD-01.\\
Course reviewer & Review feedback and scope decisions; receives the
architecture and demonstration evidence. &
The course reviewer confirms the academic scope and acceptance approach.
This is an engineering-process interaction, not a runtime electrical interface.\\
\end{tabularx}
\normalsize

\section{Proposed Acquisition Scenario}
\begin{enumerate}
\item The project owner confirms the board's input and power limits before
connecting a source.
\item The owner builds and loads the firmware using a recorded configuration.
\item The owner applies a known analog input within the approved range.
\item The system acquires ADC readings and exposes them through the selected
output mechanism.
\item The owner compares the observed readings with an independent reference
using the agreed conversion scale, tolerance, and sample timing.
\end{enumerate}
These steps describe the intended end-to-end demonstration. They are not a
test procedure or an observed result; detailed procedures and mappings will
be developed in \ChapterRef{Tests} after requirements are approved.

\section{Assumptions, Dependencies, and Exclusions}
A supervised bench setting is proposed; operating temperature, electrical
noise, available instruments, and physical access remain to be confirmed.
The project depends on access to a suitable target, its documentation,
a controllable analog source, programming tools, and an independent reference
instrument. Availability of these resources has not been verified.

Channel count, resolution, sampling schedule, latency, accuracy, filtering,
buffering, and output throughput are not fixed. No UART, USB, SPI, I2C,
RTOS, DMA, or network dependency is selected by this draft.
Custom hardware manufacture and actuators are outside the proposed initial
scope described in \ChapterRef{Introduction}; the reviewer can revise these
limits before requirements baseline.

\section{Fault and Boundary Questions}
An input outside the supported range, lost reference or ground, unstable power,
ADC timeout, missed sample, or unavailable output consumer can invalidate a
reading or damage a target. Their applicability and responses require decisions,
not assumptions. TBD-02, TBD-03, TBD-04, and TBD-05 in
\hyperref[Table:WeekOneTBDs]{the decision register} capture the required
electrical limits, timing and integrity criteria, output behavior, and
startup/fault policy. No safe operating range or fault recovery claim is made
before hardware evidence is available.

\section{Week 2 Candidate Scenario Elaboration}
\label{Section:WeekTwoScenarioContext}
The owner requested potentiometer-to-ADC0804 8-bit conversion followed by
FT245R capture, plus separate oscillator-driven 74xx counter and paired-FTDI
8-bit stream experiments on November 12, 2026. The requirements chapter
records these as candidate scenarios with explicit hardware, timing,
electrical, and protocol decisions in
\hyperref[Table:RequirementParameters]{the parameter register}.
This supersedes the absence of named candidate parts in the Week 1 draft;
it does not establish reviewed wiring, hardware availability, electrical
limits, or implementation. The counter and paired-path experiments are not
approved replacements for the lecture's single sensor acquisition scenario.

\section{Context Review Exit}
The context is ready for requirements development when the owner and reviewer
have agreed the scope, participants, boundary, and the disposition of any TBD
that affects a binding requirement. Unresolved questions remain visible with
owners and decision gates in \ChapterRef{DevelopmentPlan}.
